\section{相似性原理与像的局部结构}
\label{sec:similarity-principle-local-image}
\stfhIndex{I127}{fu^luo^er^huo^fei^er^sa^la^meng}{弗洛尔--霍费尔--萨拉蒙（Floer–Hofer–Salamon）}
本节引入一个有用的局部相似性原理，把带全实边界条件的 $J$-全纯曲线方程联系到边界条件完全用复变量表述的经典全纯曲线。借此可把 $J$-全纯曲线的许多局部研究与估计约化为经典全纯曲线的相应结论。弗洛尔、霍费尔与萨拉蒙 (FHS95) 曾把这一原理用于内部性质，本书作者 (Oh97a) 则把它用于边界性质。这里集中于边界情形，因为内部情形可包含在边界情形之中。

记 $\mathcal J(2n)$ 为满足 $I^2=-\mathrm{id}$ 的自同构
$I:\mathbb R^{2n}\to\mathbb R^{2n}$ 所成集合；称每个这样的 $I$ 为
$\mathbb R^{2n}$ 上的线性复结构。
\begin{definition}\label{def:ch8-linear-totally-real}
\stfhIndex{I047}{00-symbol-047}{$\mathcal T R(I;n)$}
设 $I:\mathbb R^{2n}\to\mathbb R^{2n}$ 为满足 $I^2=-\mathrm{id}$ 的自同构。若
$V\cap IV=0$\footnote{译者注：原书此处印作“$V$ 横截于 $JV$”，其中 $J$ 为 $I$ 之误植；译文按等价条件 $V\cap IV=0$ 表述。}，称 $n$ 维子空间 $V\subset\mathbb R^{2n}$ 关于 $I$ \emph{全实}。记这种子空间的集合为
$\mathcal{TR}(I;n)\subset Gr(2n,n)$，其中 $Gr(2n,n)$ 是
$\mathbb R^{2n}$ 中 $n$ 维子空间的集合。
\end{definition}
每个 $\mathcal{TR}(I;n)$ 都是 $Gr(2n,n)$ 的开子集，并且
\[
 \mathcal{TR}(n):=\bigcup_{I\in\mathcal J(2n)}\mathcal{TR}(I;n)
 \longrightarrow\mathcal J(2n)
\]
构成纤维丛。

令 $D_\epsilon(0)$ 为边界 $\partial D_\epsilon(0)=(-\epsilon,\epsilon)
\subset\mathbb R\subset\mathbb C$ 的开闭半圆盘，并以 $z=(x,y)$ 或 $z=x+iy$ 表示标准坐标。设
$J:z\mapsto J(z)$ 属于 $W^{1,p}(D_\epsilon,\mathcal J(2n))$，
$C:z\mapsto C(z)$ 属于 $L^p(D_\epsilon,L_\mathbb R(\mathbb R^{2n}))$，其中
$L_\mathbb R(\mathbb R^{2n})$ 是 $\mathbb R^{2n}$ 的自同态所成的集合，并假设
$p>2$。设
$\tau:(-\epsilon,\epsilon)=\partial D_\epsilon(0)\to
\mathcal{TR}(J(x);n)$ 是拉回丛 $J^*\mathcal{TR}(n)$ 的截面。由迹定理（例如见 (Rud73)），
$J|_{(-\epsilon,\epsilon)}\in W^{1-1/p,p}$；要求 $\tau$ 作为到
$Gr(2n,n)$ 的映射也属于 $W^{1-1/p,p}$，故其边界限制连续。对三元组
$((J;\tau),C)$，考虑
\begin{equation}
\begin{cases}\xi_x+J(z)\xi_y+C(z)\xi(z)=0,\\ \xi(x,0)\in\tau(x).
\end{cases}\label{eq:ch8-linear-rh-problem}
\end{equation}
稍微滥用记号，以 $i$ 同时表示复数单位和
$\mathbb C^n\cong\mathbb R^{2n}$ 上相应的标准复结构，并以
$\bar\partial$ 表示标准柯西--黎曼算子
$\partial/\partial x+i\partial/\partial y$。类似地，还以
$\bar\partial_J$ 表示线性算子
\[
 \xi\longmapsto\frac{\partial\xi}{\partial x}
 +J\frac{\partial\xi}{\partial y}.
\]
本节的主要结果是下述边值问题的相似性原理。该结果在 (Oh97a) 的定理 2.1
中得到证明，是弗洛尔、霍费尔与萨拉蒙对内部情形之结果 (FHS95) 的边界版本。
\begin{theorem}[相似性原理]\label{thm:ch8-similarity-principle}
\stfhIndex{I247}{xiang^si^xing^yuan^li}{相似性原理（similarity principle）}
\stfhIndex{I238}{li^man^xi^er^bo^te^wen^ti}{黎曼--希尔伯特问题（Riemann–Hilbert problem）}
设 $((J;\tau),C)$ 如上，$p\in(2,\infty)$，且
$\xi\in W^{1,p}(D_\epsilon(0),\mathbb R^{2n})$ 为上述问题的解。则存在
$0<\delta<\epsilon$、处处可逆的
$\Phi\in W^{1,p}(D_\delta,L_\mathbb R(\mathbb R^{2n}))$ 及
$\sigma\in C^\infty(D_\delta,\mathbb R^{2n})$，使
$\xi=\Phi\sigma$、$J\Phi=\Phi i$；若把 $\sigma$ 看作从 $D_\delta(0)$ 到 $\mathbb C^n$ 的映射，则它满足经典黎曼--希尔伯特问题
\begin{equation}\bar\partial\sigma=0,\qquad\sigma(x,0)\in\mathbb R^n\subset\mathbb C^n.\label{eq:ch8-classical-rh-problem}\end{equation}
\end{theorem}
\begin{proof}
先取
$\Psi\in W^{1,p}(D_\epsilon,GL(2n,\mathbb R))$，使
\[
 \begin{cases}
 J(z)\Psi(z)=\Psi(z)i,&z\in D_\epsilon(0),\\
 \Psi(x,0)\mathbb R^n=\tau(x),&x\in\partial D_\epsilon(0).
 \end{cases}
\]
由于边界区间可缩，选择 $\tau$ 的全局实标架
$\{\Psi_1(x,0),\ldots,\Psi_n(x,0)\}$，它们关于 $J(x)$ 复线性无关；规定
$ie_j$ 映到 $J(x,0)\Psi_j(x,0)$，并把各向量作为列，得到
$\Psi(x,0)\in GL(2n,\mathbb R)$，满足 $J(x,0)\Psi(x,0)=\Psi(x,0)i$。
因 $\tau\in W^{1-1/p,p}$，边界上的 $\Psi$ 也属该空间；再延拓到内部，使
$J\Psi=\Psi i$ 且 $\Psi\in W^{1,p}$。
\begin{exercise}\label{ex:ch8-extend-adapted-frame}证明此断言。\end{exercise}
接着由方程 $\xi(z)=\Psi(z)\eta(z)$ 定义
$\eta\in W^{1,p}(D_\epsilon(0),\mathbb R^{2n})$。把 $\eta$ 看成到
$\mathbb C^n$ 的映射，则边界条件为
\[
 \eta(x,0)\in\mathbb R^n,\qquad x\in(-\epsilon,\epsilon),
\]
而方程为
\[
 0=\bar\partial_J\xi+C(z)\xi
 =\frac{\partial\xi}{\partial x}
 +J(z)\frac{\partial\xi}{\partial y}(z)+C(z)\xi(z).
\]
这里
\[
 \frac{\partial\xi}{\partial x}(z)
 =\psi(z)\frac{\partial\eta}{\partial x}(z)
 +\frac{\partial\psi}{\partial x}(z)\eta(z),
\]
并且
\begin{align*}
 J(z)\frac{\partial\xi}{\partial y}(z)
 &=J(z)\Psi(z)\frac{\partial\eta}{\partial y}(z)
   +J(z)\frac{\partial\Psi}{\partial y}(z)\eta(z)\\
 &=\Psi(z)i\frac{\partial\eta}{\partial y}(z)
   +J(z)\frac{\partial\Psi}{\partial y}(z)\eta(z).
\end{align*}
把这些代回上一方程，便得到
\[
 \begin{cases}
 \bar\partial\eta+\widetilde C\eta=0,\\
 \eta(x,0)\in\mathbb R^n,
 \end{cases}
\]
其中
\[
 \widetilde C(z)=\Psi^{-1}(z)
 \left(\psi_x(z)+J(z)\Psi_y(z)\right)+\Psi^{-1}(z)C(z)\Psi(z).
\]
因 $C\in L^p$、$J,\Psi\in W^{1,p}$，故 $\widetilde C$ 仍属于 $L^p$。
但它不是复线性的，即未必与乘 $i$ 可交换。现在要找\emph{依赖于} $\eta$ 的复线性矩阵
$B(z)\in GL(2n,\mathbb R)$；把它看作 $GL(n,\mathbb C)$ 中的元素，使
$\widetilde C(z)\eta(z)=B(z)\eta(z)$。把
$\widetilde C\in L^p(D_\epsilon(0),GL(2n,\mathbb R))$ 分解为
\[
 \widetilde C=\widetilde C_\ell+\widetilde C_a,\qquad
 \widetilde C_\ell=\tfrac12(\widetilde C-i\widetilde C i),\quad
 \widetilde C_a=\tfrac12(\widetilde C+i\widetilde C i).
\]
记 $\Gamma:\mathbb C^n\to\mathbb C^n$ 为复共轭。定义复线性映射
\[
 D(z)=
 \begin{cases}
 |\eta(z)\rangle\langle\overline{\eta(z)}|/|\eta(z)|^2,&\eta(z)\ne0,\\
 0,&\eta(z)=0,
 \end{cases}
\]
其中 $|\eta(z)\rangle\langle\overline{\eta(z)}|$ 是狄拉克的 ket 算子记号，
对于列向量 $\eta(z)$，它表示
$(\eta(z))^t\overline{\eta(z)}$。于是
$D(z)\eta(z)=\overline{\eta(z)}=\Gamma(\eta(z))$。由于
$\eta\in W^{1,p}$ 故连续，零点集合可测；由定义 $|D(z)|\leq1$，并且
$D\in L^\infty(D_\epsilon,GL(n,\mathbb C))$。于是
\[
 \bar\partial\eta+B\eta=0,\qquad
 \eta(x,0)\in\mathbb R^n,\qquad
 B=\widetilde C_\ell+\widetilde C_a\overline D,
\]
且 $B$ 复线性。
这使我们能够证明下述引理。
\begin{lemma}\label{lem:ch8-rh-trace-fredholm}
令
\[
 \operatorname{Tr}:W^{1,p}((D_\epsilon(0),\partial D_\epsilon(0)),
 (\mathbb C^n,\mathbb R^n))
 \longrightarrow W^{1-1/p,p}(\partial D_\epsilon(0),\mathbb R^n)
\]
为迹映射。则线性映射
\[
 (\bar\partial,\operatorname{Tr}):
 W^{1,p}((D_\epsilon(0),\partial D_\epsilon(0)),(\mathbb C^n,\mathbb R^n))
 \longrightarrow L^p(D_\epsilon(0),\mathbb C^n)
 \times W^{1-1/p,p}(\partial D_\epsilon(0),\mathbb R^n)
\]
满射，且弗雷德霍姆指标为 $n$。
\end{lemma}
\begin{exercise}\label{ex:ch8-prove-rh-fredholm}证明此引理。\end{exercise}
取截断函数 $\chi:\mathbb H\to[0,1]$，使
\[
 \chi(z)=
 \begin{cases}
 1,&z\in D_{1/2}(0),\\
 0,&z\notin D_1(0),
 \end{cases}
\]
并令 $\chi_\delta(z)=\chi(z/\delta)$。则
$\|\chi_\delta B\eta\|_p\leq C\|\eta\|_{1,p}\|\chi_\delta B\|_p
\leq C(\delta)\|\eta\|_{1,p}$，其中 $C(\delta)\to0$。所以充分小 $\delta$ 时，
令 $B_\delta=\chi_\delta B$，则
\[
 (\bar\partial+B_\delta,\operatorname{Tr}):
 W^{1,p}((D_\epsilon(0),\partial D_\epsilon(0)),(\mathbb C^n,\mathbb R^n))
 \longrightarrow L^p(D_\epsilon(0),\mathbb C^n)
 \times W^{1-1/p,p}(\partial D_\epsilon(0),\mathbb R^n)
\]
仍由满射性的开性及指标在连续变形下不变而满射且指标为 $n$。令基本解
\[
\Theta_\delta\in W^{1,p}((D_\epsilon,\partial D_\epsilon),
(GL(n,\mathbb C),GL(n,\mathbb R)))
\]
满足
\begin{equation}
\begin{cases}
\bar\partial\Theta+B_\delta\Theta=0,&\text{于 }D_\delta(0),\\
\Theta(x,0)=Id\in GL(n,\mathbb R),&\text{于 }\partial D_\delta(0).
\end{cases}
\label{eq:ch8-fundamental-solution}
\end{equation}
充分小的 $0<\delta<\epsilon$ 下该解存在且唯一。定义
$\sigma:B_\delta(0)\to\mathbb C^n$ 使
$\eta=\Theta_\delta\sigma$，直接用 \eqref{eq:ch8-fundamental-solution} 计算得
$\sigma$ 满足所需方程，从而完成证明。
\end{proof}

显然，内部相似性原理已包含于定理 8.8.2。下面给出其证明关于
$J$-全纯曲线内部奇点结构的一个直接推论。
\begin{corollary}\label{cor:ch8-critical-point-expansion}
设 $J$ 是 $(M,\omega)$ 的任意相容殆复结构，并考虑非恒定
$J$-全纯映射 $u:(\Sigma,j)\to(M,J)$。若 $du(z_0)=0$，则存在
$k\in\mathbb N$ 及 $z_0$ 处的复坐标 $z$，使
$j^{k-1}u(z_0)=0$、$j^ku(z_0)\ne0$ 且
$u(z)=\vec v z^k+O(|z|^{k+1})$，其中 $\vec v\in T_{u(z_0)}M$；特别地，
存在 $z_0$ 的开邻域 $V\subset\Sigma$，使所有 $z\in V\setminus\{z_0\}$
都满足 $du(z)\ne0$。
\end{corollary}
受此推论启发，引入下述记号。
\begin{definition}\label{def:ch8-ramification-order-principal-jet}
对于非恒定 $J$-全纯映射 $u:(\Sigma,j)\to(M,J)$，称上述推论中的整数
$k$ 为 $u$ 在 $z_0$ 的分歧阶，称 $j^ku(z_0)$ 为 $u$ 的主 jet。
\end{definition}
由该推论还知，任意 $J$-全纯映射的主 jet 是全纯的，即该 jet 属于
$J^{(k,0)}_{z_0}(M)$。分歧型与临界点主 jet 的详细结构见 (Oh11a)。边界点的相应结论需要进一步研究边界行为，不过上述边界相似性原理仍然适用。

下面陈述相似性原理的几个直接推论。
\begin{proposition}\label{prop:ch8-boundary-critical-isolated}
设 $J_0$ 为 $\mathbb R^{2n}$ 上与 $(x,0)$ 无关的固定殆复结构，$R$ 为关于
$J_0$ 的全实子流形；设
$J:D_\epsilon(0)\times M\to\mathcal J(2n)$ 光滑且
\begin{equation}J(z,p)=J_0(p),\qquad z=(x,0)\in\partial D_\epsilon(0).\label{eq:ch8-boundary-independent-j}\end{equation}
若 $w:(D_\epsilon,\partial D_\epsilon)\to(\mathbb R^{2n},R)$ 光滑并满足
\begin{equation}w_x+J(z,w)w_y=0,\qquad w(x,0)\in R,\label{eq:ch8-boundary-nonlinear-cr}\end{equation}
则存在 $\delta>0$，使 $dw(z)\ne0$ 对 $0<|z|<\delta$ 成立。
\end{proposition}
\begin{proof}
若 $dw(0)\ne0$，无须证明。设 $dw(0)=0$，令 $\xi=w_x$；对给定
柯西--黎曼方程关于 $x$ 微分，得到
\[
 \begin{cases}
 \xi_x+J(z,w(z))\xi_y+\dfrac{\partial J(z,w(z))}{\partial x}w_y=0,\\
 \xi(x,0)\in T_{w(x,0)}R.
 \end{cases}
\]
又由原方程 $w_y=J(z,w(z))w_x=J(z,w(z))\xi$，故若写
\[
 C(z):=\frac{\partial J(z,w(z))}{\partial x}J(z,w(z)),
\]
则 $\xi$ 满足定理 8.8.2 的方程。因而存在
$\Phi\in C^\infty(D_\delta,GL(\mathbb R^{2n}))$ 与全纯 $\sigma$，使
$\xi=w_x=\Phi\sigma$ 且 $\sigma(0)=0$。对全纯映射，零点要么孤立，要么在
$0$ 的邻域恒为零；于是存在 $\delta>0$，使 $0<|z|<\delta$ 时
$\sigma(z)\ne0$，也即 $w_x(z)\ne0$。这就完成了证明。
\end{proof}

\begin{theorem}\label{thm:ch8-boundary-unique-continuation}
设 $J,R$ 如命题 8.8.8，并设
$w_1,w_2:(D_\epsilon(0),\partial D_\epsilon(0))\to(\mathbb R^{2n},R)$
是满足 \eqref{eq:ch8-boundary-nonlinear-cr} 的光滑映射。若它们在 $0$ 的无穷 jet 相同，则
$w_1\equiv w_2$。
\end{theorem}
\begin{proof}
不失一般性，设 $w(0,0)=0\in R\subset\mathbb R^{2n}$，且
$T_0R=\mathbb R^n\subset\mathbb C^n\cong\mathbb R^{2n}$。取边界拉平局部微分同胚
$\psi:U\to\mathbb R^{2n}$，其像是 $0$ 的开邻域，满足
$\psi(\mathbb R^n\cap U)\subset R$ 且 $\psi(0)=0$。则
\eqref{eq:ch8-boundary-nonlinear-cr} 等价于
\[
 \widetilde w_x+\psi^*J_z\widetilde w_y=0,\qquad
 \widetilde w(x,0)\in\mathbb R^n,
\]
其中 $\widetilde w=\psi^{-1}\circ w$，$J_z(w)=J(z,w)$，且
$\psi^*J_z(w)=T\psi^{-1}J(z,w)T\psi$。

因 $w_1,w_2$ 的无穷 jet 相同，$\widetilde w_1,\widetilde w_2$ 也如此。令
$\xi(z)=\widetilde w_2(z)-\widetilde w_1(z)$。直接计算表明，$\xi$ 满足
\eqref{eq:ch8-boundary-nonlinear-cr} 型方程，其中
\[
 J(z)=\psi^*J_z(\widetilde w_2(z)),
\]
而 $C(z)$ 由
\[
 C(z)\xi:=\left(\int_0^1D(\psi^*J_z)
 (s\widetilde w_1+(1-s)\widetilde w_2)\,ds\right)\xi\cdot
 \frac{\partial\widetilde w_1}{\partial y}
\]
确定。并且 $j^\infty\xi(0,0)=0$。由定理 8.8.2，
$\xi=\Phi\sigma$，其中 $\Phi(z)\in GL(\mathbb R^{2n})$ 而 $\sigma$ 全纯。
因 $\Phi(0)$ 可逆，$j^\infty\sigma(0,0)=0$，故全纯性蕴含
$\sigma=0$，进而 $\xi=0$。这证明了 $\widetilde w_1=\widetilde w_2$，从而
$w_1=w_2$ 在 $0$ 附近成立。

现取 $z_0\in D_\epsilon(0)$，考虑线段 $t\in[0,1]\mapsto tz_0$，并令
$K=\{t\in[0,1]\mid w_1(tz_0)=w_2(tz_0)\}$。将证明 $K=[0,1]$。
$K$ 显然闭，且上述局部讨论给出某个 $\delta>0$ 使 $K\supset[0,\delta]$。
令 $0<t_0\leq1$ 并定义
\[
 t_0=\sup\{t_1\mid[0,t_1]\subset K\}.
\]
我们断言 $t_0=1$。由定义，对所有满足 $0\leq t\leq t_0$ 的 $t$ 都有
$w_1(tz_0)\equiv w_2(tz_0)$。反设 $t_0<1$。在 $t_0z_0$ 处应用上述讨论，得到一个全纯映射 $\sigma$，它在
$0\leq t\leq t_0$ 的所有 $tz_0$ 上为零。于是 $\sigma$ 在 $t_0z_0$ 的一个小邻域中恒为零，特别地对充分小的
$\delta>0$，$\sigma(tz_0)$ 在 $t\in(t_0-\delta,t_0+\delta)\subset[0,1]$ 上为零。这进而蕴含
$w_1(tz_0)\equiv w_2(tz_0)$ 对 $t\in(t_0-\delta,t_0+\delta)$ 成立，从而在
$[0,t_0+\delta)$ 上成立；这与 $t_0$ 的定义矛盾。这就完成了证明。
\end{proof}
更多边界奇点结构及伪全纯圆盘像沿边界的延拓，见 (Oh97a)。
